% !TEX root = ../../Main.tex
\section{Statistical Screening and the Rejected Candidates}
\label{Section:StatisticalScreening}

Every published model had to pass a test against a null built from the model itself. The model's own sequence of positions is rotated in time. The rotated sequence holds the same positions for the same durations, but at the wrong moments. Its regime score is therefore the score that the model's holding structure earns by luck alone. The published model must beat that score.

The null must be computed for each model and cannot be a fixed number. Across the seven pairs it goes from $48.21$ on GBP/USD to $69.82$ on NZD/USD. A fixed threshold of, say, $55$ would be below chance on one pair and well above it on another. It would accept noise in one market and reject skill in another. This was found by trying a fixed threshold first.

The screening gave a result that changes how the whole campaign should be read. \textbf{On every one of the seven pairs, the candidate with the highest return failed the test.} \Cref{Tables:BetaTraps} lists the strongest cases. A portfolio built by ranking candidates on return would have chosen models that return roughly $+250\%$ on average. Almost all of that return is market exposure, not policy.

% !TEX root = ../Main.tex
\begin{table}[htb!]
\centering
\footnotesize
\caption{Candidates rejected by the screening of \Cref{Section:StatisticalScreening}. On every pair, the candidate with the highest return failed. Ranking on return alone would have chosen these models over the ones published. They show roughly five times the return on paper, but almost none of the measurable skill.}
\label{Tables:BetaTraps}
\begin{tabularx}{\textwidth}{@{}lrrX@{}}
\toprule
\textbf{Pair} & \textbf{Return} & \textbf{$p$-value} & \textbf{Why it was rejected} \\
\midrule
AUD/USD & $+246.8\%$ & $0.146$ & Failed the screen \\
AUD/USD & $+142.6\%$ & $0.987$ & Published during the campaign, then withdrawn when the test was run \\
EUR/USD & $+141.6\%$ & --- & One-sided: long for $86\%$ of the window \\
GBP/USD & $+259.8\%$ & $0.163$ & Failed the screen \\
NZD/USD & $+253.9\%$ & $0.258$ & Failed the screen \\
USD/CHF & $+463.9\%$ & $0.973$ & Positions timed no better than the same positions shuffled; followed the franc trend \\
USD/CHF & $+341.3\%$ & $0.865$ & As above \\
\bottomrule
\end{tabularx}
\end{table}

The USD/CHF case is the clearest. A candidate that returned $+463.94\%$, more than nine times the published model, has a permutation $p$-value of $0.973$. Its positions are timed no better than the same positions shuffled. Its return came from following the trend of the franc. The published model returns a ninth of that, and it is the one with measurable skill.

Two cases are reported because they are uncomfortable. An AUD/USD candidate that returned $+142.55\%$ was published during the campaign. It was withdrawn when its $p$-value came out at $0.987$. On EUR/USD, the model in place at the time won on four of six metrics, including return, Sharpe, Sortino and drawdown. It was still replaced, because it lost on the one measure that decides whether a result is real.

This screening has limits, and they should be stated. Each published model was chosen from between $128$ and $256$ candidates, so a $p$-value near $0.05$ is close to what that many draws give by chance. Under a strict multiple-testing correction, only USD/CAD and GBP/USD, both at $0.0015$, would clearly pass. The permutation test is therefore used here as a screen that rejects bad models, which it clearly did. It is not used as a claim that the remaining models are statistically significant.